\section{Background and Related Work}
\label{sec:related-work}

This paper's claim is narrow and specific: a production concern relocated
to an execution engine, and enforced there by default, is a different kind
of guarantee than the same concern addressed by a technique an agent
author must choose to apply. We situate that claim against four bodies of
prior work---declarative agent authoring, context-management technique,
secure agent execution, and benchmarking methodology---and close with the
gap each leaves open.

\subsection{Agents as code, and the move toward declarative definition}
\label{subsec:rw-declarative}

The dominant construction pattern in the mainstream agent frameworks we
evaluate throughout this paper (\S\ref{subsec:method-frameworks}) is
agents-as-code: an agent is a Python (or similar) program that
imperatively wires together a language model, a set of tools, and a
control loop. Recent work has begun to push back against this pattern by
proposing that an agent be \emph{defined} declaratively rather than
programmed. Oracle's Open Agent Specification~\cite{amini2025} proposes a
unified, portable representation for an agent's capabilities, tools, and
behavior, intended to let an agent definition be authored once and
executed across runtimes. Auton~\cite{cao2026} similarly separates an
agent's declarative specification from the framework that executes it,
targeting portability and reduced boilerplate across agentic
applications. Daunis~\cite{daunis2025} proposes a declarative language
purpose-built for describing and orchestrating LLM-powered agent
workflows, replacing hand-written orchestration code with a workflow
description the runtime interprets.

Colmena shares this family's premise---that an agent should be data
interpreted by an engine, not a program the author writes
(\S\ref{subsec:declarative})---but the family's stated goal is uniformly
\emph{authoring}: a common, portable surface for specifying what an agent
is and does, so that the same definition can move between runtimes or
be authored with less code. None of these proposals treats the
declarative surface as a place to bind production properties---context
lifecycle, credential handling, retry semantics---as engine-enforced
behavior of the node types themselves; the specification describes an
agent's structure and capabilities, and what happens inside a node
executing that structure remains outside its scope.

\subsection{Context-management technique}
\label{subsec:rw-context}

A substantial body of technique addresses the cost and scaling of the
model's context window, the concern of \S\ref{sec:context}. Provider-level
prompt and context caching~\cite{anthropic_promptcaching,
google_contextcaching} lets a framework re-send a previously seen prefix
at a reduced marginal price; this lowers the \emph{price} charged for
re-transmitted context but leaves the underlying token \emph{volume}, and
therefore the size of the payload the framework constructs and the
provider processes, unchanged---a distinction this paper holds fixed by
construction (\S\ref{subsec:method-price}). Retrieval-augmented
generation~\cite{lewis2020} addresses a different axis, replacing a static
context with material retrieved on demand from an external store,
reducing what must be included up front at the cost of a retrieval step
and an indexing pipeline the application must build and maintain.
Conversation summarization is a further, widely used technique: a
framework periodically compresses prior turns into a shorter summary
in place of the verbatim transcript. MemGPT~\cite{packer2023} generalizes
the pattern into an operating-system-style memory hierarchy, paging
context between an in-context working set and an external store as a
session grows. Dynamic tool loading~\cite{gansun2025} applies the same
retrieval principle to tool schemas rather than conversation history,
loading only the tool definitions relevant to a given step instead of
holding every tool's schema resident for the whole session.

Each of these techniques is effective, and each is, in the framework
that offers it, an \emph{option}: something an agent author selects,
configures, and wires into their own application code, on a
tool-by-tool or turn-by-turn basis, if they know the hazard and choose to
address it. None is a default property of the frameworks we evaluate in
\S\ref{sec:context}; our steelman arm demonstrates this directly, closing
the token gap only after we hand-wrote the equivalent of one such
technique into the framework's own tool code
(\S\ref{subsec:context-closure}).

\subsection{Secure agent execution and secret handling}
\label{subsec:rw-security}

A production agent frequently mediates a credential---an API key, a
session token, a one-time password---between a user and a downstream
tool or endpoint, and the language model sits, by construction in most
frameworks, on the path that credential travels. Because the model's
context is exactly the artifact that is logged, cached, and
re-transmitted on every subsequent turn, any secret admitted into it is
thereby admitted into the transcript, the provider's request history, and
every downstream store the transcript touches. This hazard is
increasingly recognized in practitioner and platform guidance around
agent deployment, though it has not, to our knowledge, been treated as a
problem to be closed by an execution engine's default behavior rather
than by an application author's discipline; \S\ref{sec:credentials}
evaluates exactly this gap, isolating the credential's inbound and
outbound paths as two enforcement points an engine can guarantee rather
than two hazards an author must remember to mitigate.

\subsection{Benchmarking methodology: self-reported versus authoritative measurement}
\label{subsec:rw-benchmarking}

Reported comparisons between agent frameworks typically rely on figures a
framework computes about itself---an internal token counter, a
self-logged cost estimate---which is convenient but couples the
measurement to the same code path being evaluated: a framework's internal
accounting can be stale relative to the provider's tokenizer, can omit a
call path the framework does not instrument, or can simply be wrong, and
none of these failure modes is visible from inside the framework being
measured. This paper adopts a different basis, developed in full in
\S\ref{sec:methodology}: every framework under test is driven through a
single shared proxy in front of the model provider, and every token count
and dollar figure this paper reports is read from the proxy's own record
of the request and response that actually crossed the wire, not from any
framework's self-report. The measurement is therefore provider-authoritative
and uniform across every framework compared, including Colmena.

\subsection{The gap}
\label{subsec:rw-gap}

Read together, these strands leave one gap open. The declarative-agent
proposals of \S\ref{subsec:rw-declarative} relocate how an agent is
\emph{authored}---a portable specification in place of hand-written
orchestration code---but leave what happens inside execution as
unconstrained as it was before: nothing in Open Agent
Spec~\cite{amini2025}, Auton~\cite{cao2026}, or Daunis's
language~\cite{daunis2025} binds context lifecycle, credential handling,
or retry semantics to the engine that interprets the specification. The
context-management and security techniques of
\S\ref{subsec:rw-context}--\S\ref{subsec:rw-security} are real and
individually effective, but each is a technique a developer must
recognize the need for, select, and wire into their own application
code, framework by framework and tool by tool; none is a property the
runtime supplies unconditionally to every agent that passes through it.
In neither strand does a production \emph{property}---a guarantee that
holds for every agent, by construction, whether or not its author knew to
ask for it---get relocated to the execution engine as an enforced
invariant. That relocation, and the consequence that a concern moved
across the boundary in \S\ref{sec:model} need not be re-authored,
correctly, in every agent that uses it, is this paper's contribution, and
it is what \S\ref{sec:context}--\S\ref{sec:hardening} measure.
